\documentclass[11pt]{article}
\input{packages.tex}
\usepackage{soul}
\DeclareMathOperator{\sign}{sign}
\DeclareMathOperator{\Imm}{Im}
\DeclareMathOperator{\rg}{rg}
\DeclareMathOperator{\Mat}{Mat}

% Настройка колонтитулов
\pagestyle{fancy}
\fancyhf{} % Очистить все поля
\fancyhead[L]{А зачем нам \ \st{что-то доказывать} кузнец? Не, нам кузнец не нужен.}
\fancyhead[R]{\href{https://t.me/lost_luna}{Автор}}
\fancyfoot[C]{\thepage}               % Справа внизу

\setlength{\headheight}{13.6pt}

\hypersetup{
    colorlinks=true,
    linkcolor=blue,
    citecolor=blue,
    filecolor=blue,
    urlcolor=blue,
}

\begin{document}

\begin{titlepage}
    \centering
    \vspace*{\fill}  % заполняет пространство сверху

    {\Huge Лекция 3. Ранг фундаментальной матрицы. Линейные пространства.
    \par}

    \vspace*{\fill}  % заполняет пространство снизу
\end{titlepage}

\setcounter{page}{1}

\section{Ранг фундаментальной матрицы}

\begin{theorem}
    Пусть $A$ --- матрица размера $m \times n$ ($m$ уравнений и $n$ переменных). Тогда: 
    \[\forall F: \ AF = 0 \ \hookrightarrow \ \rg F \leq n - \rg A\]
\end{theorem}

\begin{proof}
    Элементарными преобразованиями приведем матрицу $A$ к упрощённому виду:
    \begin{figure}[H]
    \centering
    \includegraphics[width=0.54\linewidth]{pictures/Приведение матрицы к упрощенному виду.png}
    \end{figure} \leavevmode
    \\
    Так как производились только ЭО со строками и столбцами, то матрицу $A'$ можно представить в виде \[A' = P \cdot A \cdot Q, \quad 
    \text{где } P \text{ и } Q \text{ --- соответствующие элементарные матрицы}\]
    Пусть матрица $F' = Q^{-1} \cdot F$. Так как $Q^{-1}$ --- невырожденная матрица, то по лемме, сформулированной и доказанной в предыдущей лекции, получаем:
    \[\rg F' = \rg F\]
    Рассмотрим следующее произведение:
    \[A' \cdot F' = P \cdot A \cdot Q 
    \cdot Q^{-1} \cdot F = P \cdot A \cdot F = O\]
    При этом, расписав произведение $A' \cdot F'$ другим способом, получим:
    \begin{figure}[H]
    \centering
    \includegraphics[width=0.54\linewidth]{pictures/Ранг фундаментальной матрицы.png}
    \end{figure} \leavevmode
    \\
    Тривиально видеть, что в матрице $F'$ первые $\rg A$ строк нулевые (иначе не получилась бы 
    нулевая матрица). Тогда, так как в матрице $F$ всего $n$ строк, то, учитывая равенство рангов $F$ и $F'$, получим:
    \[\rg F = \rg F' \leq n - \rg A\]
\end{proof}

\begin{remark}
    Пусть $n - \rg A \geq 1$. Тогда можно выбрать $F$ так, что 
    её ранг равен в точности $n - \rg A$.
\end{remark}

\begin{proof}
    Для доказательства сего чуда предъявим $n - \rg A$ ЛНЗ в совокупности столбцов высоты $n$, из которых и соберётся искомая матрица $F$.
    Элементарными операциями только со строками (и перестановкой столбцов,
    если требуется) приведем матрицу $A$ к следующему виду:
    \begin{figure}[H]
    \centering
    \includegraphics[width=0.5\linewidth]{pictures/Упрощенный вид номер не знаю какой.png}
    \end{figure} \leavevmode
    \\
    В преобразованной матрице $\rg A$ базисных переменных и $n - \rg A$
    параметрических. Пусть следующие столбцы и образуют искомую фундаментальную матрицу $F$:
    \begin{figure}[H]
    \centering
    \includegraphics[width=0.54\linewidth]{pictures/Столбцы фундаментальной матрицы.png}
    \end{figure} \leavevmode
    \\
    Эти столбцы ЛНЗ, так как ни одна $1$ не может быть получена линейной
    комбинацией нулей. Следовательно, \[ \rg F = n - \rg A\]
\end{proof}

\begin{statement}
    Любое решение системы представимо в виде линейной комбинации столбцов
    фундаментальной матрицы.
\end{statement}

\begin{proof}
    Допишем к $F$ произвольный столбец решений: $\big( F \ \vdots \ x \big)$.
    Ранг матрицы $F$ не мог поменяться по определению фундаментальной матрицы. 
    Тогда, по теореме Кронекера-Капелли получаем,
    что $x$ является линейной комбинацией столбцов фундаментальной матрицы.
\end{proof}

\begin{definition}
    ФСР (фундаментальная система решений) --- множество решений, которые
    являются столбцами фундаментальной матрицы.
\end{definition} \leavevmode
\\
Следствием утверждения выше является то, что $x = F \cdot c$ --- общее решение системы уравнений $A \cdot x = 0$, где $c$ --- некоторый столбец коэффициентов. \\
Рассмотрим же теперь неоднородную систему $A x = b$. Пусть мы нашли некоторое её решение
решение $h_1$. Поскольку любые два решения неоднородной системы отличаются
друг от друга на решение однородной системы $A x = 0$ (это было сформулировано и доказано в предыдущей лекции), то, пристально посмотрев
на умозаключение выше, получим, что общее решение неоднородной системы $A x = b$ есть:
\[ x = h_1 + F \cdot c\]

\begin{remark}
    Как же найти сие дивное $h_1$?
\end{remark} \leavevmode
\\
Элементарными преобразованиями строк (и перестановкой столбцов, если требуется),
приведем расширенную матрицу системы к следующему виду (здесь мы конечно же уже подразумеваем, что система совместна):
\begin{figure}[H]
\centering
\includegraphics[width=0.54\linewidth]{pictures/Дивное решение.png}
\end{figure} \leavevmode
\\
По этому виду вдумчивый читатель сразу смекнет, что решением будет столбец

\[
\begin{array}{c}
h_1 = \left( \begin{array}{c}
 \\ \overline{b}' \\ \\ 0 \\ \vdots \\ 0
\end{array} \right)
\begin{array}{l}
\left. \begin{array}{c} 
 \\  \\ \\
\end{array} \right\} \rg A \\[5pt]
\left. \begin{array}{c} 
 \\ \\ \\
\end{array} \right\} n - \rg A
\end{array}
\end{array}
\]

\section{Линейные пространства}

\begin{definition}
    Множество $\mathcal{L}$ называется \textbf{линейным пространством}, а его элементы
    называются \textbf{векторами}, если на этом множестве заданы операции сложения элементов 
    \["+" \colon \mathcal{L} \times \mathcal{L} \to \mathcal{L}\]
    и умножения элементов на вещественное число
    \["\cdot" \colon \mathbb{R} \times \mathcal{L} \to \mathcal{L}\] 
    относительно которых это множество замкнуто, то есть 
    \[\forall x, y \in \mathcal{L} \hookrightarrow x + y \in \mathcal{L}, \quad
    \forall x \in \mathcal{L} \; \; \forall \alpha \in \mathbb{R} \hookrightarrow \alpha \cdot x \in \mathcal{L}\]
    и которые удовлетворяют следующим аксиомам:
    \begin{enumerate}
        \item \[ \forall x, y \in \mathcal{L} \hookrightarrow x + y = y + x\]
        \item \[ \forall x, y, z \in \mathcal{L} \hookrightarrow (x + y) + z = x + (y + z) \]
        \item \[ \exists 0 \in \mathcal{L} : \forall x \in \mathcal{L} \hookrightarrow 0 + x = x\] 
        \item \[ \forall x \in \mathcal{L} \; \; \exists (-x) \in \mathcal{L} : x + (-x) = 0\]
        \item \[ \forall \alpha, \beta \in \mathbb{R} \; \; \forall x \in \mathcal{L}
        \hookrightarrow (\alpha \beta) x = \alpha (\beta x) \]
        \item \[ \forall x \in \mathcal{L} \hookrightarrow 1 \cdot x = x\]
        \item \[ \forall \alpha, \beta \in \mathbb{R} \; \; \forall x \in \mathcal{L} 
        \hookrightarrow (\alpha + \beta) x = \alpha x + \beta x \]
        \item \[ \forall \alpha \in \mathbb{R} \; \; \forall x, y \in \mathcal{L} 
        \hookrightarrow \alpha (x + y) = \alpha x + \alpha y \]
    \end{enumerate}
\end{definition}

\begin{remark}
    Множество $\mathcal{L}$ непусто, так как $\exists 0 \in \mathcal{L}$.
\end{remark}

\begin{definition}
    Пространство, состоящее только из $0$, с корректно определенными операциями,
    называется нулевым.
\end{definition}

\begin{remark} Определить разность векторов в линейном пространстве можно следующим образом:
    \[ x - y = x + (-1) \cdot y \]
\end{remark}

\newtheorem*{consequence}{Следствие}

\begin{consequence}
    \[(-1) \cdot y = -y\]
\end{consequence}

\begin{consequence}
    Нулевой элемент единственен.
\end{consequence}

\begin{proof}
    Пусть есть 2 нулевых элемента $0_1$ и $0_2$. Тогда
    \[ 0_1 + 0_2 = 0_2 + 0_1 = 0_2 = 0_1\]
\end{proof}

\begin{consequence}
    Противоположный элемент единственен.
\end{consequence}

\begin{proof}
    Пусть есть 2 противоположных элемента к $a \in \mathcal{L}: \ -a_1$ и $-a_2$.
    Тогда \[a - a_1 = 0 = a - a_2 \implies a_2 = a_1\]
\end{proof}

\begin{example}
    Линейное пространство из действительных положительных чисел
    (как говорили великие, возьмем абстракцию над операциями 
    умножения и сложения в линейном пространстве):
    \[ a + b = a \cdot b\] 
    \[ \alpha \cdot a = a^\alpha\]
\end{example}

\begin{statement}
    ФСР для системы $Ax = 0$ образует линейное пространство.
\end{statement}

\begin{remark}
    Понятия линейной комбинации, линейной зависимости и линейной 
    независимости можно распространить на линейные пространства. При 
    этом свойства и теоремы, которые не использовали геометрические свойства 
    пространства, также справедливы.
\end{remark}

\begin{definition}
    \textbf{Базисом} в линейном пространстве $\mathcal{L}$ назовём упорядоченную \textbf{конечную} систему 
    векторов, удовлетворяющую следующим условиям:
    \begin{enumerate}
        \item Эта система ЛНЗ
        \item Любой вектор из $\mathcal{L}$ раскладывается по этой системе
    \end{enumerate}
\end{definition}

\begin{remark}
    Коэффициенты разложения любого вектора из $\mathcal{L}$ по базису единственны и они называются координатами
    вектора в этом базисе.
\end{remark}

\begin{statement}
    Пусть в пространстве $\mathcal{L}$ есть базис. Тогда при сложении векторов их координаты
    складываются. Аналогично, при умножении вектора на вещественное число 
    координаты вектора умножаются на это число.
\end{statement}

\begin{consequence}
    Линейная комбинация векторов $\Longleftrightarrow$ линейная комбинация 
    их координатных столбцов \[ \text{ЛЗ (ЛНЗ) векторов} \quad 
    \Longleftrightarrow \quad \text{ЛЗ (ЛНЗ) их координатных столбцов}\]
\end{consequence}

\begin{theorem}
    Если в линейном пространстве $\mathcal{L}$ есть базис из $n$ векторов, то 
    любой набор из  $m > n$ векторов линейно зависим.
\end{theorem}

\begin{proof}
    Пусть $\{ e_1, e_2, \dots, e_n\}$ --- базис в $\mathcal{L}$. Пусть 
    $\{ f_1, f_2, \dots, f_m\}$, где $m > n$, --- набор векторов.
    Рассмотрим координатные столбцы векторов $\{ f_1, f_2, \dots, f_m\}$
    в базисе $\{ e_1, e_2, \dots, e_n\}$:
    \[ \left(
        \begin{array}{c}
            \\ f_1 \\ \\
        \end{array}
    \right), \ \left(
        \begin{array}{c}
            \\ f_2 \\ \\
        \end{array}
    \right), \dots , \left(
        \begin{array}{c}
            \\ f_m \\ \\            
        \end{array}
    \right) \left. \begin{array}{l}
        \\ \\ \\
    \end{array}\right\} \  n
    \]
    Составив из них матрицу, получаем, что в ней столбцов больше, чем строк, а значит по теореме, сформулированной в первой лекции, имеем, что все эти столбцы ЛЗ, а тогда по нашему блестящему следствию
    вектора также ЛЗ.
\end{proof}

\begin{consequence}
    У любых двух базисов в $\mathcal{L}$ количество векторов одинаково.
\end{consequence}

\begin{proof}
    Предположим, что это не так, то есть в каком-то базисе векторов больше.
    Тогда, по предыдущей теореме получим, что эти вектора линейно зависимы,
    а значит они не могут быть базисом. Получаем противоречие, что и требовалось.
\end{proof}

\begin{definition}
    \textbf{Размерностью} линейного пространства $\mathcal{L}$ будем назвать количество векторов в базисе $\mathcal{L}$. Размерность 
    пространства $\mathcal{L}$ обозначается $\dim \mathcal{L}$.
\end{definition}

\begin{remark}
    По договорённости размерность нулевого пространства равна $0$.
\end{remark}

\begin{example}
    Рассмотрим линейное пространство, состоящее из матриц размера $m \times n$. Легко видеть, что:
    \[\dim \Mat_{m \times n}(\mathbb{R}) = m \cdot n\]
    В самом деле, базисом в этом пространстве можно выбрать $m \times n$ матриц вида 
    \begin{figure}[H]
    \centering
    \includegraphics[width=0.4\linewidth]{pictures/Базисные матрицы.png}
    \end{figure}
\end{example}

\begin{example}
    Рассмотрим линейное пространство, состоящее из многочленов степени
    не больше $n$. Тогда очевидно, что его размерность равна $n + 1$, а базисом можно взять следующий набор:
    \[\{1, t, t^2, \dots, t^n\}\]
\end{example}

\begin{definition}
    Пусть линейное пространство $\mathcal{L}$ таково, что $\forall m \in \mathbb{N}$ в нём существует набор 
    из $m$ линейно независимых векторов. Тогда пространство $\mathcal{L}$ назовём \textbf{бесконечномерным}.
\end{definition}

\begin{remark}
    Пространство $C([a, b])$ бесконечномерно.
\end{remark}

\begin{figure}[H]
    \centering
    
\includegraphics[width=0.8\linewidth]{pictures/ну я.jpg}
    \caption{Автор после лекции представляет, как будет все это техать}
    \end{figure}
\end{document}